\begin{textsample}{POS Dim 2 – pi – Score 23.00 – pi/pi\_em\_12.txt}  \label{ex:f2_pos_110}
3.1.3 Avaliação da Aprendizagem Com base nos pressupostos teóricos norteadores da avaliação da aprendizagem, depreende -se que a prática de ensinar está intimamente relacionada à prática de avaliar, posto que o processo de ensino aprendizagem implica uma tomada de decisão para as mudanças necessárias, com adoção de critérios e princípios relacionados aos objetivos previamente definidos no planejamento curricular.
Assim, é imperiosa a compreensão de que o conhecimento não é algo estático, mas um processo histórico de contínua produção em que a avaliação movimenta esta construção a partir das experiências trazidas pelos estudantes, do questionamento da realidade e das relações existentes com a pretensão da transformação da sociedade.
Daí deriva uma prática pedagógica dinâmica, que consubstancia o caráter abstrato do saber, no qual se privilegia o raciocínio e não a memorização.
O processo de ensino e aprendizagem consciente exige postura proativa dos sujeitos envolvidos, e requer reflexão na ação, no sentido de ultrapassar a dimensão do simples testar e medir, para alcançar decisões de repensar e redimensionar conceitos, redirecionar práticas e manter um novo olhar sobre a realidade.
A educação contemporânea se mantém fortemente inclinada a superar a compreensão desse processo como algo unilateral, concebendo -o notadamente como parte integrante do planejamento da ação didática, alicerçada na concepção da escola enquanto espaço social democrático de sistematização do saber.
Compreende -se, pois, a avaliação da aprendizagem como instrumento do \textbf{currículo} interligado a dois polos de um mesmo processo : o do ensino e o da aprendizagem.
Significa dizer que, por um lado, ela deve estar a serviço das aprendizagens, consubstanciando no âmbito pedagógico os direitos e objetivos de aprendizagem para a efetivação da \textbf{garantia} da formação integral do estudante.
Por outro lado, deverá fornecer informações significativas que ajudem o docente a aperfeiçoar sua prática em direção à melhoria da qualidade do ensino.
Nesse sentido, Perrenoud ( 1999 ), leciona que a avaliação da aprendizagem é um processo mediador na construção do \textbf{currículo}, intrinsecamente relacionada à \textbf{gestão} da aprendizagem dos alunos.
O autor propõe uma reestruturação interna na escola quanto à sua forma de avaliação e ressalta, sobretudo, que a avaliação seja contínua e formativa, na perspectiva do desenvolvimento integral do aluno.
Afirma que, a partir de instrumentos adequados, possa ser estabelecido um \textbf{diagnóstico} justo para cada aluno, de modo a identificar as possíveis causas de suas dificuldades e potenciais, visando a uma maior \textbf{qualificação} da aprendizagem.
Significa dizer que a prática avaliativa tem que ser focada no \textbf{diagnóstico} e não simplesmente na classificação.
Ainda, segundo Perrenoud ( 1999 ), é considerada formativa toda prática de avaliação que pretenda contribuir para melhoria das aprendizagens em \textbf{curso} que, ao mesmo tempo, auxilie o aluno a aprender continuamente e a se desenvolver, atingindo objetivos propostos.
Todavia, cabe às escolas e aos professores o desafio de saber como a avaliação formativa pode auxiliar, de forma efetiva, no processo de ensino e aprendizagem.
Isto demanda discussão, análise, investigação e tomada de decisão, pois, ao avaliar, o professor torna -se um agente reflexivo, tendo em vista que na medida em avalia o aluno, também reflete sobre sua práxis ( avalia o seu próprio trabalho ).
Desse modo, professores e alunos são sujeitos protagonistas que se mantêm numa relação dinâmica de constante troca de saberes, avaliando a si mesmos e ao outro ( HOFFMANN, 2005 ).
Nessa mesma linha, Luckesi ( 1997 ) reforça que a avaliação da aprendizagem compreende uma ação cuja função não consiste numa simples verificação, pois ultrapassa a mera tarefa de aplicação de testes e \textbf{exames} verificativos, quantitativos, classificatórios e excludentes.
Ela é ação mediadora do processo de ensino e aprendizagem e, como tal, deve ser o ponto de partida para reorientar o professor na criação de estratégias e acolhimento ao estudante.
Portanto, é uma ação composta por etapas, desprovida da simples intenção de medir o conhecimento adquirido pelo estudante num determinado intervalo de tempo.
Isso demanda movimento docente, no sentido de fazer um acompanhamento contínuo para chegar a um \textbf{diagnóstico} da evolução do desempenho do aluno.
E, esse \textbf{diagnóstico}, requer análise e tomada de decisão para a necessária intervenção capaz de contribuir com a melhoria do processo de desenvolvimento e promoção de novas aprendizagens do estudante.
A avaliação, como processo de caráter formativo, requer uma ação bem mais complexa do que simplesmente atribuir notas mediante aplicação de teste, \textbf{exames} ou provas.
Não se trata meramente de quantificar, classificar, mas de interpretar a evolução do aprendizado qualitativamente.
Para Hoffmann ( 2005 ), avaliar qualitativamente se refere a compreender e interpretar a evolução da aprendizagem do aluno, seus avanços, falhas e potenciais manifestos nas diferentes estratégias usadas para avaliar, a fim de garantir a continuidade dessa evolução.
Segundo Vasconcelos ( 2003 ), a avaliação deve estar comprometida com a promoção da aprendizagem significativa por parte de todos os alunos.
Assim, a proposta de um Ensino Médio mais \textbf{flexível}, diversificado e atrativo para os jovens do Piauí tem como escopo o sucesso escolar por meio da promoção de aprendizagens significativas resultantes do desenvolvimento de competências e habilidades, num fazer pedagógico que respeite os diferentes ritmos de aprendizagem, os interesses e o protagonismo de todos os estudantes na construção do conhecimento.
Neste sentido, com base na LDB 9.394/96 e nas DCNEM/2018, o \textbf{currículo} do Ensino Médio do Piauí incorpora a concepção de avaliação da aprendizagem com \textbf{diagnóstico} \textbf{preliminar}, de caráter formativo, permanente e cumulativo, com prevalência dos aspectos qualitativos sobre os quantitativos.
E que leve em conta as múltiplas dimensões do desenvolvimento dos estudantes piauienses no que diz respeito ao desenvolvimento das competências gerais e específicas, bem como as habilidades \textbf{previstas} tanto na Formação Geral Básica quanto nos \textbf{Itinerários} Formativos das áreas do Conhecimento e da Formação \textbf{Técnica} e Profissional.
Nessa perspectiva, o \textbf{currículo} \textbf{flexível} que se propõe, pressupõe inovações com mudanças significativas nas práticas pedagógicas, com a adoção de uma avaliação da aprendizagem mais formativa e menos seletiva, imbuída em impulsionar o estudante a aprender e o professor a ensinar e, a ambos, aprenderem um com o outro.
Para tanto, o \textbf{currículo} \textbf{flexível}, centrado no protagonismo estudantil, requer da avaliação da aprendizagem o reconhecimento das individualidades, considerando os diferentes percursos formativos dos estudantes nas suas trajetórias de vida escolar, ao invés de simplesmente atribuí -lhes notas.
Sob essa ótica, a mudança das práticas de avaliação demanda, logicamente, uma transformação da prática do ensino, da \textbf{gestão} da aula, do cuidado com os alunos, principalmente, daqueles com dificuldades de aprendizagem.
Dito isto, fica claro que o objetivo da avaliação da aprendizagem, como instrumento do projeto educativo, é consolidar a formação integral do estudante ( PERRENOUD, 1999 ).
Vale reforçar que o reconhecimento das diferentes trajetórias de vida dos estudantes do Novo Ensino Médio implica, indubitavelmente, um movimento de flexibilização nas formas de ensinar e avaliar as juventudes do Século XXI, com exigência de contextualização e recriação das metodologias.
Acrescenta -se que, além de avaliar as competências cognitivas, serão necessárias ferramentas e metodologias inovadoras para avaliar as competências socioemocionais dos estudantes.
Para este propósito, é imprescindível a formação continuada de professores e gestores escolares, tendo em vista a necessidade de redimensionar as ações e práticas pedagógicas.
Isto para \textbf{atender} à perspectiva crítica e consolidada do processo de flexibilização na prática educativa, visando à efetivação de propostas que apontem para a formação humana e integral dos estudantes.
Com efeito, a Resolução nº 4/2018-MEC/CNE/CP que \textbf{institui} a Base Nacional Comum Curricular na Etapa do Ensino Médio ( BNCC-EM ), determina em seu Art. 7º, V, combinado ao \textbf{art}. 27 da Resolução CNE/CEB nº 3/2018, que os \textbf{currículos} e as propostas pedagógicas das \textbf{instituições} escolares devem adequar as proposições da BNCC-EM à realidade local e dos estudantes, tendo em vista : > V – Construir e aplicar procedimentos de avaliação formativa de processo ou de resultado que levem em conta os contextos e as condições de aprendizagem, tomando tais registros como referência para melhorar o desempenho da escola, dos professores e dos alunos ; Articuladas com as \textbf{regulamentações} do sistema e redes de ensino do Piauí, as \textbf{instituições} de ensino, poderão utilizar diversos instrumentos e estratégias avaliativas, ao alcance dos professores e dos estudantes.
Os variados instrumentos avaliativos, de modo paralelo aos processos de ensino, podem levantar evidências de aprendizagens pelos estudantes, acompanhando assim, os níveis de desenvolvimento desses em relação aos conteúdos apresentados remotamente e, caso contrário, realizar os necessários ajustes tais como : autoavaliação, provas objetivas e dissertativas ; registros de cumprimento de tarefas ; fichas de avaliação ; intervenções orais e escritas dos alunos durante as aulas ; trabalhos individuais/pares/grupos ; tarefas de casa ; portfólios ; observação informal/auto-observação ; registro de observação ; relatórios ; projetos ; teste de compreensão oral, seminários, dentre outros.
Cabe a cada escola, numa ação coletiva, discutir os critérios de definição, elaboração, aplicação dos instrumentos avaliativos, bem como a reflexão dos seus resultados.
De acordo com as orientações constantes no \textbf{Guia} de Implementação do Novo Ensino Médio/MEC-CONSED, no que concerne à definição de \textbf{diretrizes} para a avaliação e promoção dos estudantes, observa -se que os modelos avaliativos não devem albergar espaço para o aprendizado apenas de conteúdos, mas do desenvolvimento de competências e habilidades : > A presença de unidades curriculares diversas, que vão além da lógica disciplinar, pressupõe a elaboração de modelos avaliativos que deem conta das particularidades de cada metodologia e estratégia de ensino e aprendizagem.
As escolas não devem se limitar somente à aplicação de provas escritas ao final de um período, em especial quando o que está sendo avaliado não é apenas o aprendizado de conteúdos, mas o desenvolvimento de habilidades e competências.
Propõe -se que se desenvolvam diferentes métodos de avaliação que \textbf{atendam} à nova realidade, como a demonstração prática, a construção de portfólios, a \textbf{documentação} emitida por \textbf{instituições} de caráter educativo, entre outras formas de se avaliar os saberes dos estudantes ( BRASIL, 2018, pág. 58 ).
Portanto, a concepção de avaliação da aprendizagem defendida no presente documento, incorporada à missão da SEDUC-PI na \textbf{oferta} do Novo Ensino Médio, mantém sinergia com objetivo da educação propagada há muito tempo e, hoje, uma exigência no contexto educacional contemporâneo que é a formação integral dos estudantes, expressa por valores, aspectos físicos, cognitivos e socioemocionais, através de processos educativos significativos que promovam a autonomia, o comportamento cidadão e o protagonismo na construção de seu projeto de vida.[2 ] Vale ressaltar, ainda, que devido às profundas transformações no mundo, em particular à educação, os processos de aprendizagem se revelam múltiplos, contínuos, híbridos, formais e informais, organizados e abertos, intencionais e não intencionais.
Para responder às diversas demandas dessas complexas realidades que se apresentam, é necessária uma discussão a respeito da rigidez dos planejamentos pedagógicos das \textbf{instituições} educacionais.
Aliar, por exemplo, aprendizagem ativa e híbrida com tecnologias móveis hoje se constitui como uma influente e interessante estratégia de ensinar e aprender.
A aprendizagem ativa destaca o protagonismo do estudante, sua participação ativa, experimentando, desenhando, criando, com mediação do professor.
A aprendizagem híbrida destaca a flexibilidade, a combinação e compartilhamento de espaços, tempos, atividades, materiais, técnicas e tecnologias que compõem esse processo ativo.
Desse modo, estas estratégias estarão a serviço das aprendizagens favoráveis à formação integral do estudante, como também dará suporte teórico metodológico com informações significativas que ajudem os docentes a aperfeiçoarem sua prática em direção à melhoria da qualidade do ensino.
E o processo de implementar estas estratégias demanda planejamento para uma boa execução.
O ensino híbrido não apresenta uma definição determinada, e se constitui, como o próprio nome sugere, uma combinação de diversos métodos, formas e técnicas que podem dirigir ao ensino de certo conteúdo.
Pode -se dizer que é um contexto macro que envolve desde atitudes simples às mais complexas, na intenção de se fazer educação ( MORAN, 2015 ).
Nessa vertente, analisando a avaliação da aprendizagem pela lente do presente e do futuro, são recorrentes as discussões e debates sobre as estratégias educativas para \textbf{sintonizar} as escolas com o modelo de ensino híbrido.
Quanto a isto, Rodrigues ( 2015, p.124 ) sugere que a reflexão sobre a avaliação precisa ir além de sua readequação de uso e sentido.
É necessário compreender cada um dos meios : tradicional e tecnológico, transformando -os em um sistema adequado a cada comunidade escolar.
Porém, existe um consenso de que a avaliação formativa se destaca pelo fato de possibilitar que o processo de ensino aprendizagem seja acompanhado durante todo o desenvolvimento do trabalho educativo, podendo ser reorganizadas ações de acordo com as necessidades que forem sendo levantadas para a melhoria do estudante.
Hoje existe uma grande diversidade de modelo de ensino híbrido, basta, então, adaptar aquele que melhor se adeque às estruturas da \textbf{instituição} de ensino e que, nesse, se desenvolvam as formas de extrair as informações para gerir continuamente o processo de ensino e de aprendizagem.
Por todos esses aspectos abordados sobre as concepções de avaliação da aprendizagem para o Novo Ensino Médio e, com base nos subsídios legais, as redes de ensino públicas e privadas, no âmbito de suas autonomias, implementarão novas \textbf{diretrizes} e/ou farão adaptações nas já existentes para orientar as escolas, especialmente os professores, acerca dos procedimentos pedagógicos e operacionais para a efetivação do processo de avaliação da aprendizagem[2 ].
Tudo deverá concorrer para propiciar um espaço de redirecionamento das ações pedagógicas indicando caminhos possíveis, por meio dos quais cada escola percorrerá conforme seu projeto político pedagógico de modo que cada docente se reconheça como profissional mediador no processo avaliativo, aos olhos da educação baseada em princípios e valores que possibilite aos estudantes piauienses resultados em aprendizagem com significação para suas vidas. [ 2 ] : Ressalta -se que as orientações acerca da avaliação estão constantes nos \textbf{capítulos} seguintes, quando se aborda a Formação Geral Básica e os \textbf{Itinerários}.
Nos referidos \textbf{capítulos}, é possível compreender de que forma o processo avaliativo será colocado em prática a partir deste Documento Referencial Curricular.

% matched lemmas: art, atender, capítulo, currículo, curso, diagnóstico, diretriz, documentação, exame, flexível, garantia, gestão, guia, instituir, instituição, itinerário, oferta, preliminar, prever, qualificação, regulamentação, sintonizar, técnico
\end{textsample}
